% Three-tier Peer-HBM/DRAM/storage research draft. Evaluation curves are analytical.
% Compile this directory independently with pdfLaTeX and BibTeX.
\documentclass[sigplan,10pt,anonymous,nonacm]{acmart}
\setkeys{acmart.cls}{balance=false}
\usepackage{subcaption}
\usepackage{tabularx}
\usepackage{placeins}
\setcopyright{none}
\settopmatter{printacmref=false,printccs=false,printfolios=true}
\acmConference[EuroSys '27]{European Conference on Computer Systems}
 {April 19--23, 2027}{Rabat, Morocco}
\acmYear{2027}
\acmDOI{}
\acmISBN{}
\title{AegisKV: Workflow-Aware KV Checkpointing\texorpdfstring{\\}{ }Across HBM, DRAM, and Shared Storage}
\setlength{\emergencystretch}{1.5em}
\renewcommand{\bibfont}{\normalfont\normalsize}
\newcounter{designalgorithm}
\captionsetup[subfigure]{font=normalsize,labelfont=normalfont,textfont=normalfont}
\AtBeginDocument{\typeout{AEGISKV-GEOMETRY: width=\the\textwidth;
 height=\the\textheight; columnsep=\the\columnsep; baseline=\the\baselineskip}}
\begin{document}
\fancypagestyle{aegispreview}{\fancyhf{}\renewcommand{\headrulewidth}{0pt}
 \renewcommand{\footrulewidth}{0pt}\fancyfoot[C]{\normalsize\thepage}}
\pagestyle{aegispreview}
\begin{abstract}
Multi-agent applications accumulate reusable context while executing
interdependent LLM requests. A worker failure can lose this state and delay
several branches at once, yet protecting every context in fast memory is
expensive. We propose AegisKV, a workflow-aware manager for KV checkpoints
across Peer HBM, shared DRAM and shared storage. A complete checkpoint
spans the three tiers, with one managed owner per object: selected layers
and hot blocks reside on a preassigned takeover DP, and lower tiers
retain the rest. AegisKV combines normal reuse with the workflow cost
of recovery, allocating request value to reusable prefixes and evaluating
requests affected by a common failure together. Sparse demand and layer
overlap guide Peer acceleration; renewable credit manages DRAM and storage.
Actual reuse renews per-Agent credit, while value updates and tier moves
preserve paid rent. Complete placement plans pass a current workflow-cost
check under bounded decision work and a migration threshold.
Versioned transitions preserve recovery validity during migration and
explicit fallback. A per-DP AIMD Adapter limits Peer capacity while the Master
chooses per-request layer depth and hot blocks without changing native
attention semantics. We connect trace-derived
workflow observations and existing reconstruction microbenchmarks to a
reproducible analytical study and a complete experimental protocol.
In 36 declared Peer-disabled scenarios, the placement model improves on
LRU in 17 and ties in 19. Matched analytical schedules illustrate Peer's
incremental overlap benefit; neither model measures a three-tier system. Historical outage timing drives a separate recovery-load model in which
a lower mean coexists with worse tail latency. Mechanism ablations expose
interactions between layer placement and overlap, and a negative result
for an earlier rental screen. Bounded repairs reduce that model's gap
but do not establish the new controller's performance. The study
identifies when contention, prediction error, and commit
lag limit the benefit. End-to-end fault injection and normal
service measurements remain necessary to establish the proposed system's
performance.
\end{abstract}

\maketitle
\thispagestyle{aegispreview}
\raggedbottom

\section{Introduction}
\label{sec:introduction}

A multi-agent application completes when its required results are ready:
a tested patch, a report, or an answer assembled from parallel searches.
An LLM request's value therefore depends on its workflow position. A
critical-branch delay postpones completion; another branch may absorb it
before a join. Application-aware serving exploits this structure during
normal execution~\cite{lin2024parrot,pan2025kvflow,ni2026topas}. Failures
raise a further question: which histories should remain quickly
recoverable when one worker interrupts several requests together?

Long contexts make that choice expensive. Agents accumulate reasoning,
code, and tool results; losing their GPU-resident KV can require substantial
reconstruction. Checkpointing avoids that work~\cite{strati2024dejavu,jayakody2026ghostserve},
and load-aware recovery distributes interrupted requests~\cite{cao2026lumen}.
Hierarchical caches provide DRAM and storage capacity~\cite{qin2025mooncake,xie2026strata},
but keeping every history in the faster tier is costly.

Recency and recovery latency alone do not determine placement. A long,
recent history may belong to a branch that tolerates storage access,
while a smaller history lies on the critical path. High slack does not
justify deleting reusable history: subsequent normal calls may repeatedly
reconstruct it. Common-worker failures also couple requests through the
DAG and shared recovery resources. Accelerating only one of two delayed
branches may leave their join unchanged. Finally, demotion is asynchronous:
freeing DRAM before its storage replacement commits invalidates protection.

AegisKV manages checkpoints across Peer HBM, shared DRAM and shared storage.
Its Tracker estimates remaining workflow time; the Master compares
complete plans under normal reuse and common-domain failures. Prefix
increments distribute request value across KV objects. Causal sparse
demand and layer overlap guide per-request Peer layer/block placement,
while DRAM and storage hold the remaining state. Objects have one managed
owner after migration. Lower-tier rent ages byte occupancy; actual Agent
reuse renews its credit, while value updates and tier moves preserve paid
rent. Peer acceleration is revalued directly. A complete-plan check
accounts for coupled victims and migration under bounded decision work.
The Router binds a compatible takeover DP at arrival; an AIMD Adapter
bounds Peer capacity and the Connector applies safe transfers.
Versioned manifests, reservations and reader leases preserve checkpoint
coverage as ownership moves between tiers.

Our proposed contributions are a workflow objective that separates Peer
acceleration, lower-tier demotion and final-copy deletion; a bounded online controller and safe transition
protocol; and reproducible analytical scenarios with a matched physical
experiment plan. The model exposes gains, ties, prediction sensitivity,
and regret against a feasible oracle. Trace-driven recovery and independent
mechanism models expose queue-tail regressions and an adverse historical
rental screen, whose bounded repairs leave a residual gap. These results
motivate and constrain the proposed
system; they do not measure its implementation. The central empirical
hypothesis is that workflow-aware protection improves finite-resource
allocation after accounting for healthy-service and migration costs.

% Background and Motivation: trace structure and matched physical KV costs.
\section{Background and Motivation}
\label{sec:background-motivation}

\subsection{Background}
\label{sec:motivation-background}

\paragraph{Multi-agent workflows.}
Coding assistants and research systems increasingly organize LLM calls
into multi-agent workflows~\cite{wu2023autogen,anthropic2025multiagent}.
An agent may have a predefined role, such as planner or reviewer, or be
created for a particular subtask. It can issue multiple LLM requests,
invoke tools, and incorporate their results before producing an output.
Sequential dependencies, parallel branches, and joins connect these
agents. A request, an agent, and the complete user workflow therefore
have different completion conditions.

\paragraph{Prefill and reusable state.}
Prefill processes a request's input tokens and produces per-layer key and
value (KV) tensors. Subsequent computation can reuse this state instead
of processing the same prefix again. Agent histories often grow as later
requests include earlier reasoning, code, and tool outputs. Reusing their
KV saves computation, but retaining every growing history consumes memory.
Hierarchical caches extend GPU memory with host DRAM and larger storage
tiers~\cite{qin2025mooncake,gao2024cachedattention}.

\paragraph{Failures at serving scale.}
A GPU or worker failure can interrupt an entire model replica, removing
both serving capacity and access to its active KV cache~\cite{cao2026lumen}.
Surviving KV copies allow recovery without reconstructing the complete
history; otherwise, retained tokens must be processed again.
KV checkpointing and load-aware recovery address these
costs~\cite{strati2024dejavu,jayakody2026ghostserve,cao2026lumen}.
We consider fail-stop failures with a compatible surviving replica,
retained token history and control state, and backup storage outside the
failed domain. Within this scope, we ask which histories deserve scarce
fast-tier capacity when their recovery affects a dependent workflow.

\begin{figure*}[t]
  \centering
  \captionsetup[subfigure]{font=normalsize,labelfont=normalfont,
    textfont=normalfont,justification=centering,singlelinecheck=false}
  \begin{subfigure}[t]{2.48in}
    \vspace{0pt}
    \begin{minipage}[c][1.65in][c]{\linewidth}
      \centering
      \includegraphics[width=\linewidth]{figures/workflow.pdf}
    \end{minipage}
    \caption{Four-branch workflow}
    \label{fig:motivation-workflow}
  \end{subfigure}\hfill
  \begin{subfigure}[t]{1.86in}
    \vspace{0pt}
    \begin{minipage}[c][1.65in][c]{\linewidth}
      \centering
      \includegraphics[width=\linewidth]{figures/workflow-gaps.pdf}
    \end{minipage}
    \caption{Workflow-maximum gaps}
    \label{fig:motivation-branch-gap}
  \end{subfigure}\hfill
  \begin{subfigure}[t]{2.52in}
    \vspace{0pt}
    \begin{minipage}[c][1.65in][c]{\linewidth}
      \centering
      \includegraphics[width=\linewidth]{figures/recovery-cost.pdf}
    \end{minipage}
    \caption{Context and Prefill cost}
    \label{fig:motivation-recompute}
  \end{subfigure}
  \caption{Workflow structure and the cost of reusable KV state.
    (a) Four trace-derived branches with schematic Main stages A0/A5.
    Bars preserve timing; dashed tails assume an all-branch join.
    (b) ECDF of maximum branch-duration gaps per workflow: 121 workflows
    with temporally inferred fan-in; these are not failure-time slack.
    (c) Measured Qwen2.5-7B-Instruct-1M Prefill on two A100 PCIe 40\,GB GPUs
    (TP2, BF16, batch one), with sparse attention above 32K.
    Reuse starts from $(C-1)$ cached tokens and computes one;
    DRAM/shared-storage I/O is serialized with computation.
    HBM/DRAM: median and min--max of three timings;
    shared storage/full prefill: one timing per length.
    Protocol: Appendix~\ref{sec:motivation-protocol}.
    K~=~1,024; M~=~1,000,000.}

  \label{fig:motivation}
  \Description{Left: a four-branch workflow with A2 completing last at
    151.6 seconds. Middle: an ECDF of workflow-maximum branch-duration
    differences with a median of 138 seconds and a 90th percentile of
    1173 seconds. Right: four measured Prefill curves for Qwen2.5-7B-Instruct-1M
    from 8K to 1,010,000 tokens: blue circles for HBM reuse, orange squares
    for serialized DRAM reuse, purple diamonds for serialized shared-storage
    reuse, and red triangles for full prefill. Both axes are logarithmic.}
\end{figure*}

\subsection{Motivation}

\paragraph{Workflow completion requires different protection priorities.}
\label{sec:motivation-workflow}
We use workflow completion time (WCT), from task release to completion of
all required results, as the application-level objective.
Figure~\ref{fig:motivation}(\subref{fig:motivation-workflow}) illustrates
why request latency alone is insufficient. A2 finishes last, 151.6\,s
after the fork, whereas A1 finishes at 60.4\,s. If the following Main stage
requires all four outputs, accelerating A1 alone does not advance the
join, while delaying A2 postpones downstream work. A branch can absorb
some additional delay without postponing the required output; we call
this tolerance \emph{slack}.

The public SemiAnalysis CC/Weka trace contains 393 filtered Claude Code
sessions~\cite{semianalysis2026cctraces}. We retain overlapping subagent
groups whose completed branches are followed by a Main request after the
last branch finishes; Main may also execute while the branches run.
This temporal fan-in proxy retains 286 groups in 121 workflows.
For each group we take the longest minus the shortest branch duration,
then retain each workflow's maximum gap.
Figure~\ref{fig:motivation}(\subref{fig:motivation-branch-gap}) gives a
median of 138\,s and a 90th percentile of 1,173\,s. These conditional
statistics establish substantial branch-duration heterogeneity. They do
not establish result-consumption edges, which the trace does not record,
or the slack distribution at a fault. The four-branch example illustrates
one retained group, not a representative sample of workflow maxima.

\paragraph{Long contexts make fast recovery state expensive.}
\label{sec:motivation-long-context}
Losing an agent's KV can turn a brief interruption into substantial
reconstruction work. Figure~\ref{fig:motivation}(\subref{fig:motivation-recompute})
compares four paths using Qwen2.5-7B-Instruct-1M~\cite{qwen2025qwen25oneM}
on two A100 PCIe 40\,GB GPUs, with TP2, BF16 and batch one.
All paths use the same SGLang 0.5.10rc0 model execution and
MInference vertical/slash attention~\cite{jiang2024minference} with dual chunk
attention (DCA) position handling; attention is sparse above 32K.
For total context $C=M+1$, full prefill starts from empty KV, while reuse
starts from the same $M$-token history in HBM, DRAM, or shared storage and
computes one token. Timing ends with all $C$ tokens' KV ready, before
vocabulary projection. HBM provides a resident-state reference; the
DRAM and shared-storage paths serialize restoration and computation.

At 1,010,000 tokens, full prefill takes 571.21\,s despite sparse attention.
HBM, DRAM and shared-storage reuse take 1.305, 2.663 and 10.156\,s,
respectively. A reusable history avoids most reconstruction work, but its
storage tier still changes readiness time: shared storage adds 7.49\,s
over DRAM in this run. The corresponding history occupies 57.92\,GB
across the two ranks. Keeping many such histories in DRAM can therefore
be expensive, while moving them to storage introduces additional delay.
These measurements characterize one fixed synthetic input and serialized
paths, with the timing counts stated in the caption. They establish the
cost difference in this configuration; storage-cache state and overlapping
I/O can change its magnitude. Appendix~\ref{sec:motivation-protocol}
gives the measurement and validation details.

\begin{figure}[t]
  \centering
  \includegraphics[width=\linewidth]{figures/recovery-exposure.pdf}
  \caption{Measured Prefill costs for
    Figure~\ref{fig:motivation}(\subref{fig:motivation-workflow}), using
    Figure~\ref{fig:motivation}(\subref{fig:motivation-recompute})'s configuration
    at trace input-size proxies plus one token.
    Bars/whiskers show median/min--max over three timings after warmup.
    Hatching denotes absorbed delay; solid fill is the calculated WCT
    increase under an all-branch join with other execution fixed.
    Each case is a single-branch counterfactual, not a fault-injected
    WCT measurement.}
  \label{fig:motivation-recovery-exposure}
  \Description{A1 through A4 have measured median Prefill costs of
    33.2, 32.6, 35.4 and 39.9 seconds.
    Applying each cost separately as an added branch pause gives WCT
    increases of 0.0, 32.6, 16.3 and 18.0 seconds.
    Hatching identifies the absorbed portion; whiskers show the minimum
    and maximum of three timings, not confidence intervals.}
\end{figure}

\paragraph{Failure position determines the workflow benefit of protection.}
Reconstruction cost and workflow position must be considered together.
Figure~\ref{fig:motivation-recovery-exposure} uses the four branches in
Figure~\ref{fig:motivation}(\subref{fig:motivation-workflow}). We measure
full prefill at each active request's input-size proxy, approximately
101K--114K tokens in decimal units, using the configuration above.
The proxies come from trace hash-block counts; the measured input is
synthetic, rather than the original Claude prompts.
At 29.5\,s after the fork, we insert each measured cost $R_i$ as an
additional pause on branch $i$, separately, while holding other execution
and queues fixed. Under
the assumed all-branch join, the trace completion times give offline
slack $\sigma_i$, and the resulting WCT increment is
$\max(0,R_i-\sigma_i)$.

A1's 33.2\,s cost fits within its 91.2\,s slack and adds no WCT
in this construction. A2 takes 32.6\,s but has zero slack,
exposing its entire cost. A3 and A4 expose 16.3 and
18.0\,s, respectively. Thus, reconstruction latency alone does
not determine the benefit of protection. This example combines physical
reconstruction timings with a conditional dependency calculation; it
does not measure failures in the original workflow or establish the
frequency of such outcomes. Common-worker failures and shared recovery
resources require joint analysis, developed in the next section.

\label{sec:motivation-goals}
These results motivate AegisKV's central objective: allocate finite
fast-tier capacity and backing storage according to their effect on
workflow completion. Retaining a usable copy avoids reconstruction;
keeping KV closer to execution can additionally avoid restore waits.
Neither saving has the same WCT value for every branch. Protection must also account for normal reuse
and the capacity and I/O it consumes. The design therefore separates the
value of fast-tier retention from the value of keeping any recoverable
copy, and evaluates both in the context of the workflow.

\section{AegisKV Design}
\label{sec:design}

AegisKV allocates finite KV storage to reduce workflow completion time
(WCT) while preserving normal cache reuse. Two observations connect
workflow-level priorities to the state needed for recovery.

\subsection{Observations}
\label{sec:design-insights}

\paragraph{1. Workflow dependencies determine recovery urgency.}
Figure~\ref{fig:motivation-recovery-exposure} shows similar reconstruction
costs with different effects on WCT: slack absorbs some branches' delays,
while a critical branch exposes its full pause. Thus request recovery
time alone is insufficient to rank protection priorities; the request's
position in the remaining workflow also matters.

\paragraph{2. Historical KV accesses concentrate on a subset of blocks.}
Figure~\ref{fig:sparse-releases} surveys long-context releases with block,
token or compressed-entry selection. Figure~\ref{fig:design-insights}
then examines native access traces: approximately 20\% of historical
blocks account for 87.64\%, 74.76\% and 56.54\% of accesses across the
three models. These results distinguish the complete historical state
from the subset frequently used by computation. The survey does not
establish a monotonic increase in sparsity, and retrospective concentration
does not establish future-query prediction accuracy.

\begin{figure}[t]
 \centering
 \includegraphics[width=\columnwidth]{figures/sparse-models.pdf}
 \caption{Selected sparse/hybrid model releases and reported context
 lengths. IDs: 1 MiniCPM4~\cite{minicpm4model};
 2 MiniCPM4.1~\cite{openbmb2025minicpm41};
 3 DeepSeek-V3.2-Exp~\cite{deepseek2025v32exp};
 4 DeepSeek-V3.2~\cite{deepseek2025v32model};
 5 GLM-5~\cite{zai2026glm5model};
 6 MiniCPM-SALA~\cite{openbmb2026sala};
 7 DeepSeek-V4 preview~\cite{deepseek2026v4model};
 8 MiniMax-M3~\cite{minimax2026m3};
 9 GLM-5.2~\cite{zai2026glm52};
 10 LongCat-Flash-Lite-Sparse~\cite{longcat2026sparse}.
 Dates are announcements, or first official repository dates for
 IDs 5, 9 and 10. Markers identify the sparse branch's selection unit,
 including SALA's block branch and V4's compressed sparse branch;
 they do not imply uniform attention across layers.}
 \label{fig:sparse-releases}
 \Description{Ten selected sparse or hybrid attention releases between
 June 2025 and July 2026, with reported contexts from 32K to 1M.
 Circles, squares and diamonds distinguish block, token and compressed
 entry selection. This is a literature survey, not a sparsity-rate curve.}
\end{figure}

\begin{figure}[t]
 \centering
 \includegraphics[width=\columnwidth]{figures/kv-access-concentration.pdf}
 \caption{Native KV access concentration for NOSA-8B~\cite{huang2026nosav2},
 NSA-8B (MiniCPM-4-8B-NSA)~\cite{fasa2026minicpm4nsa}, and
 MiniCPM4.1-8B~\cite{openbmb2025minicpm41}. We use the final 512 queries
 of matched 32,704-token prefixes from three PG19 texts. Each query
 counts a historical 64-token block once across heads. Rank the 503
 blocks within each layer/window; average access shares over all
 32 layers, then texts. Bands show text min--max; the diagonal denotes
 uniform access. Compressed/selector state is excluded. This is
 retrospective concentration, not prediction accuracy or recovery speedup.}
 \label{fig:design-insights}
 \Description{Cumulative historical block access curves for NOSA-8B,
 NSA-8B and MiniCPM4.1-8B, averaged over every layer and three texts.
 About 20 percent of blocks account for 87.64, 74.76 and 56.54 percent
 of accesses, respectively; all three curves exceed the uniform baseline.}
\end{figure}

Together, these observations guide which states to protect and which
parts to keep immediately accessible.

\subsection{System Overview}
\label{sec:system-overview}

AegisKV combines workflow priorities with a three-tier KV hierarchy
(Figure~\ref{fig:system-overview}). Peer HBM keeps selected
layers and hot blocks on another data-parallel (DP) replica for fast
recovery; shared DRAM and storage retain the remaining state for reuse
and recovery. Together, the three tiers cover each committed checkpoint,
with one managed owner per object after transfers finish. This division lets
the system retain long histories while spending scarce HBM on the state
whose faster access is most useful to workflow completion.

The Application supplies requests and workflow progress to the Controller.
\textbf{(1)} Workflow Tracker estimates how much delay each Agent can
absorb, while \textbf{(2)} Recovery-aware Router selects its execution and
takeover DPs. Using this urgency and location information, \textbf{(3)} KV
Master decides where to retain each KV object and which hot blocks merit
Peer space. Within the Engine, \textbf{(4)} Adapter adjusts the available
Peer capacity as serving demand changes. Connector executes Master's
transfer plans and reports their outcomes. If a worker fails, Router redirects its requests to the
takeover DP, where Connector combines Peer hits with lower-tier reads to
resume execution.

\begin{figure}[t]
 \centering
 \includegraphics[width=\columnwidth]{figures/system-overview.pdf}
 \caption{AegisKV architecture. Callouts (1)--(4) match the text;
 DP1 mirrors DP0's Adapter/Connector; the Master--Connector exchange is
 shown on DP1: plans flow down and transfer status returns up.
 DP0 executes; DP1 hosts selected Peer objects and takes over. Blue solid links
 carry KV, dark dashed links carry control, and red dashes denote retry;
 orange blocks are unsealed tails. The remaining checkpoint objects
 reside in shared DRAM/storage outside both DP failure boundaries.}
 \label{fig:system-overview}
 \Description{Compact Application, Controller and Engine layers with
 numbered components, symmetric DP Local/Peer HBM and shared lower tiers.}
\end{figure}

\subsection{Workflow Tracker: Remaining Time and Slack}
\label{sec:workflow-protection}

To make workflow priorities available online, Tracker estimates remaining
execution before deriving slack. It maintains a residual DAG whose nodes
are Agent segments between exposed fork/join boundaries or completion.
A segment can span several LLM calls and tool waits. Child execution
appears on child paths, not again in the parent; unknown future forks
are excluded.

A GBDT predicts segment seconds remaining, $d$, from causal features: role/model,
current input length, elapsed time, completed-call count and previous
calls' lengths, timings and context growth. Missing history is masked;
training, calibration and evaluation split whole workflows. Labels end
at the segment boundary, and future call outcomes are never features.

Algorithm~\ref{alg:workflow-tracker} takes residual DAG $G$ and features
$X$. Forward/backward critical-path passes~\cite{kelley1959cpm} compute
earliest finishes $EF$ and latest finishes $LF$ relative to the predicted
workflow finish. Their difference is nominal slack. The guard $g$ is an
upper empirical quantile of positive slack overestimation on calibration
workflows; subtracting it gives advisory slack $s$, not a per-request
guarantee. Requests in one segment share its margin.

\par\medskip\noindent\begin{minipage}{\columnwidth}
\raggedright
 \refstepcounter{designalgorithm}\label{alg:workflow-tracker}
 \textbf{Algorithm \thedesignalgorithm: Remaining time and slack}
 \par\smallskip
 \begin{tabularx}{\columnwidth}{@{}r@{\quad}X@{}}
 \toprule
 1 & $(G,X)\gets\textsc{SnapshotProgress}()$ \\
 2 & $\nu\gets\textsc{NextVersion}()$ \\
 3 & \textbf{if not} \textsc{Supported}$(G,X)$ \textbf{then} \\
 4 & \quad \textsc{PublishInvalid}$(\nu)$; \textbf{return} \\
 5 & $d\gets\max(0,\textsc{GBDT}(X));\quad g\gets\textsc{Guard}()$ \\
 6 & $EF\gets\textsc{Forward}(G,d)$ \\
 7 & $LF\gets\textsc{Backward}(G,d,\max_v EF_v)$ \\
 8 & \textbf{for} $v\in V(G)$ \textbf{do} \\
 9 & \quad $s_v\gets\max(0,LF_v-EF_v-g_v)$ \\
 10 & \textsc{Publish}$(\nu,G,d,g,s)$ \\
 \bottomrule
 \end{tabularx}
\end{minipage}\par\medskip

\textsc{Supported} checks inputs, dependencies, calibration and freshness;
invalid estimates grant no positive slack credit. Progress events and
periodic refresh publish versioned durations $d$, guards and slack to
Master, invalidating older placement costs. The DAG passes cost
$O(|V|+|E|)$; prediction quality is evaluated in
Section~\ref{sec:eval-prediction}.

\subsection{Recovery-aware Router}
\label{sec:recovery-router}

Router establishes where normal and resumed execution run. Normal
dispatch balances load and compatible prefix locality.
Figure~\ref{fig:agent-prefix-reuse} reports 97.90\% potential same-agent
token coverage and 0.37\% additional cross-agent coverage. Agent affinity
is a preference: shared prefixes can still help another Agent's first call.

\begin{figure}[t]
 \centering
 \includegraphics[width=\columnwidth]{figures/agent-prefix-reuse.pdf}
 \caption{Potential prefix reuse over 98,827 requests, 2,090 agents and
 393 CC/Weka workflows. Coverage uses earlier completed requests with
 the same recorded model and workflow. Same/other-agent coverage
 overlaps; only additional cross-agent coverage extends same-agent
 reuse. These are token-prefix opportunities, not runtime cache hits.}
 \label{fig:agent-prefix-reuse}
 \Description{Same-agent prefixes dominate aggregate reuse; cross-agent
 history can also help first subagent calls.}
\end{figure}

At arrival, Router records a versioned takeover binding to a DP with a
compatible layout outside the source failure domain. Its corresponding
tensor-parallel (TP) shards can own selected Peer objects. On confirmed
failure, Router fences the old epoch and activates that target subject
to workspace and admission constraints. A new target cannot count the
old target's Peer objects as local: rebinding retrieves reachable objects
or uses an older checkpoint/replay fallback if their sole owner was lost.
This binding gives Master concrete access paths and target load for
placement evaluation.

\subsection{KV Master: Workflow-aware Tiered Placement}
\label{sec:kv-master}

Master turns Tracker's workflow estimates and Router's takeover binding
into four decisions: which requests benefit from retention, which KV
objects carry that value, where to place them, and how to revise the
placement as demand changes.

\subsubsection{Workflow-aware Request Value}
\label{sec:placement-value}

A request's value comes from both normal reuse and failure recovery.
Retaining a prefix avoids reconstruction on a later call; faster recovery
also reduces WCT when its delay exceeds the request's guarded slack.
A large slack can absorb a recovery pause without making repeated
reconstruction free: those GPU cycles still compete with other requests.
Master therefore includes both benefits, using visible reuse history and
the current residual workflow. An existing local hit earns no additional
retention benefit.

Slack alone is insufficient to rank requests. The amount of work saved
and the capacity required also matter (Figure~\ref{fig:retention-value}).
Moreover, a source-worker failure can interrupt several branches at once;
accelerating only one of two tied branches may not advance their join.
Master ultimately compares complete placements $P$ and $P'$ through
\begin{equation}
 \Delta J=\sum_e p_e\sum_w[T_w(P',e)-T_w(P,e)].
 \label{eq:value}
\end{equation}
Here $T_w$ is predicted remaining WCT. The scenario table covers normal
continuation and source-domain faults at specified offsets, with
probabilities summing to one. Both plans use the same queues, finite
workspace and service profiles. Transfer interference enters request
execution times; the residual DAG propagates them through joins using
the maximum. Tracker's guard is applied once per interrupted segment.
Invalid workflow estimates fall back to request delays without positive
slack credit; invalid resource profiles pause optional moves.

\begin{figure}[t]
 \centering
 \includegraphics[width=\columnwidth]{figures/retention-value.pdf}
 \caption{Slack alone does not rank retention value. In this isolated-fault
 example, A uses two capacity units to reduce recovery from 4 to 2\,s;
 B uses one to reduce 18 to 8\,s. A at slack 0 saves 1\,s/unit, whereas
 B at slack 6 saves 10\,s/unit. Normal-use value is zero. These constructed
 profiles motivate cost/size comparison, not a policy-speedup claim.}
 \label{fig:retention-value}
 \Description{A higher-slack object can save more workflow time per byte.}
\end{figure}

\subsubsection{Prefix and Block Value}
\label{sec:prefix-value}

Request value must reach the prefix that makes reuse possible, rather
than accumulate on the newest block. For each Agent, Master starts from
a valid fallback and extends the recoverable prefix in token order.
Each increment includes the new token extent across all layers, TP shards
and required selector state. Its grant combines normal-use savings and
recovery delay removed beyond guarded slack, using the same horizon and
scenarios as Equation~\eqref{eq:value}. The previous restore path remains
an option, so increments are nonnegative. Restore times of 12, 7 and 5\,s
at successive frontiers, with slack 6\,s, yield recovery-value increments
of 5 and 1\,s before probability weighting. The final block does not
receive the entire 6\,s benefit.

Each increment is distributed over its newly required physical objects
in proportion to bytes. Distinct Agents contribute separate shares to a
shared prefix; repeated references by one Agent do not create new grants.
The physical object is charged once. Thus earlier prefixes can accumulate
value from more users, while later state retains its incremental benefit.
This allocation is a capacity-screening rule, not a claim that each layer
has independent recovery value. Grants are recomputed relative to the
available restore path; removing a dependency changes the value of the
remaining plan. Their local sum does not replace the joint WCT check.

Retention and immediate accessibility also need different priorities.
Keeping a recoverable prefix requires its full dependencies. Deciding
which of those objects deserve Peer HBM instead depends on when recovery
will access them and whether a lower-tier transfer would cause a wait.

\subsubsection{Three-Tier Placement and Peer Acceleration}
\label{sec:sparse-placement}

AegisKV uses exclusive placement across its managed tiers. A sealed
layer/token-block/TP object has one owning tier: Peer HBM, shared DRAM or
shared storage. A checkpoint is complete when the union of these objects
covers its prefix, including native auxiliary state. This managed state
is separate from the execution GPU's working KV. Moving an object to Peer
releases its lower-tier copy after transfer completes; capacity is not
spent on a permanent duplicate (Figure~\ref{fig:sparse-recovery}).

For request $r$, $L_r$ is the Peer prefix depth and $K_{r,l}$ counts the
selected token blocks in layer $l\le L_r$. Mandatory native dependencies
come first; other blocks are ranked by selection frequency over completed
queries, counted once per physical block per query across heads. Ties
favor recency, then block ID. A budget below the mandatory set is
infeasible. Selection units, engine blocks and kernel tiles need not
coincide; 64 tokens is a measurement unit here, not a universal requirement.
Objects outside the selected sets remain in DRAM or storage, and demand
misses retrieve their actual owners without changing native selections.

\begin{figure}[t]
 \centering
 \includegraphics[width=\columnwidth]{figures/sparse-recovery.pdf}
 \caption{One checkpoint distributed across three exclusive tiers.
 Each cell denotes a layer/block group, including its TP objects.
 Peer covers $L=2$ with $K_1=2$ and $K_2=1$; DRAM and storage hold the
 remaining cells. Layer 1 demands B0/B1/B2: B0 is local, B1 needs H2D,
 and B2 needs a storage read followed by H2D. The dashed tail is not yet
 committed. Demotion transfers ownership down the hierarchy; it cannot
 discard Peer bytes on the assumption of an existing lower copy.}
 \label{fig:sparse-recovery}
 \Description{A three-layer four-block ownership map with Peer, DRAM
 and storage owners, plus the three corresponding recovery access paths.}
\end{figure}

The last complete checkpoint covers $M$ tokens; the acknowledged tokens
beyond it form replay tail $N$. Saved history avoids reconstruction, while
computation on the tail can hide transfers for later layers. Peer earns
space only for the additional WCT reduction from avoiding exposed
lower-tier reads/H2D, relative to placing the same objects below it.
This benefit depends on $M$, $N$, hot-set coverage and recovery queues.
Peer uses direct revaluation, with no rental clock.

Figure~\ref{fig:design-pipeline} illustrates this distinction. Moving hot
sets from DRAM to Peer removes about 24\,ms of exposed wait while retaining
the same 1,024\,MiB of unique managed KV. The useful seven-layer depth is
specific to that history, tail and profile.

\begin{figure}[t]
 \centering
 \includegraphics[width=\columnwidth]{figures/recovery-pipeline.pdf}
 \caption{Analytical layered recovery at equal unique KV coverage.
 MiniCPM4-8B has 32 layers, 32K history and a 64-token tail. Replay is
 calibrated at 2.870\,ms/layer; assumed H2D bandwidth is 6.25\,GB/s
 (5.369\,ms/full layer). Trace-derived mean preload and preload-plus-miss
 fractions are 0.6584/0.7748. (a)--(c) hold 1,024\,MiB in DRAM; (d) moves
 147.5\,MiB to Peer in the first seven layers, leaving 876.5\,MiB in DRAM.
 Storage is inactive in this comparison. Numbers are tail completion
 and exposed wait; cold KV follows replay. Placement is already complete,
 and the ideal preemptible stream omits movement, finite-workspace and
 foreground-queue costs. These are analytical schedules, not end-to-end
 recovery measurements.}
 \label{fig:design-pipeline}
 \Description{Four recovery schedules with identical replay computation
 and unique KV coverage, reducing exposed wait from 172 to 83, 44 and
 20 milliseconds through overlap, sparse preloading and Peer placement.}
\end{figure}

Master enumerates bounded block-count/fraction templates and depths from
zero to model depth, retaining the current placement as an option. Each
fault scenario uses current $M$ and its predicted $N$, including commit
lag; unsupported sparse profiles use full demand. Only computation that
must still execute hides transfer. Neither $L$ nor $K$ must grow with
every token: block identities can change at fixed size, and $L=0$ is valid.

Lower-tier placement and Peer selection are evaluated together. New cold
state may enter DRAM or storage directly. Peer demotion moves objects to
DRAM; DRAM demotion moves them to storage. A full destination may require
a cascade before Peer memory becomes reusable. Final storage eviction
must publish an older valid checkpoint or input/log replay fallback for
every affected manifest. A shared object has one owner even when Agents
have different takeover targets: only its host DP gets a local Peer hit;
other access paths and the cost of moving it enter the plan evaluation.
The protected source failure must leave all advertised owners reachable;
loss of a sole owner requires fallback, not an assumed hidden replica.

\subsubsection{Renewable Credit and Online Management}
\label{sec:placement-plan}

Repeated calls and shared prefixes create demand that outlives one request.
A one-time score misses new users, while renewing credit at every
checkpoint can preserve stale state indefinitely. Master uses renewable
credit in DRAM and storage, inspired by Landlord and
GreedyDual-Size~\cite{young2002landlord,cao1997greedydual}, while revaluing
Peer acceleration directly as described above.

For lower pool $j$, object $x$ has size $b_x$, Agent grants $g_{a,x,j}$ and
paid amounts $d_{a,x,j}$. With cumulative pressure rent $\rho_j$, its
remaining credit and eviction key are
\begin{equation}
 C_{x,j}=\sum_a[g_{a,x,j}-d_{a,x,j}]^+,
 \qquad H_{x,j}=\rho_j+C_{x,j}/b_x.
 \label{eq:rental-credit}
\end{equation}
Each pool starts at zero. A rent increase $\delta$ charges every committed
resident object $\min(C_{x,j},b_x\delta)$, proportionally across its positive
Agent shares. Pinned objects also age; removability is a separate property.
Rent measures capacity competition rather than elapsed time.

Actual reuse renews only the using Agent's share; a use ID prevents repeated
renewal within that use. Value changes refresh grants while preserving
paid amounts. Checkpoint growth and tier migration are not reuse.
Per-pool paid amounts survive tier round trips, but accrue only while
resident in that pool; a move through Peer cannot erase earlier rent.
A new use resets that Agent's paid amounts. Tool waits preserve forecast
interest; expired interest removes grants. Figure~\ref{fig:renewable-credit}
illustrates why renewal differs from revaluation.

\begin{figure}[t]
 \centering
 \includegraphics[width=\columnwidth]{figures/renewable-credit.pdf}
 \caption{Lower-tier credit for one shared, unit-size DRAM object.
 Initial Agent grants are 3 and 2. Rent of 2 leaves shares 1.8 and 1.2.
 Updating A's grant to 4 preserves its paid amount 1.2, leaving 2.8;
 an actual new use by A then renews its share to 4. B stays at 1.2.
 Credit screens capacity competition; it is neither measured WCT nor
 duplicated capacity. Peer residency uses no rent.}
 \label{fig:renewable-credit}
 \Description{A lower-tier ledger separating pressure aging, value-only
 refresh and actual-use renewal for two Agents sharing one object.}
\end{figure}

\paragraph{Admission and victim search.}
For DRAM/storage, Master refreshes $q\ge2$ low-key objects and $q$ rotating
alternatives within a common work budget. Victim bundles must release
actual bytes without breaking manifests. Releasing lower-tier objects
$Y$ raises that pool's trial clock to at least $\max_{x\in Y}H_{x,j}$.
Incoming bytes must retain positive credit after the same trial charge;
existing objects receive no duplicate grant. Peer victim candidates use
current recovery loss per byte, subject to the same legal-plan check.
Bounded bundles can also advance tied workflow branches together.

Ordinary competition respects running requests' logical Peer pins.
Request end/tool handoff unpins; self hot-set changes and required quota
relief can override logical pins, never physical read/copy leases.
Every plan reserves destination and temporary space, recursively including
lower-tier victims. Optional plans require positive incoming lower-tier
credit and $\Delta J<-\epsilon$. Mandatory capacity relief bypasses the
positive-credit gate and chooses the least costly examined legal plan.
If no safe cascade fits, admission waits; shrinking a quota does not
instantly free occupied HBM.

\paragraph{Closing the loop.}
Prefill offers newly sealed state; decode triggers review every configured,
block-aligned $Z$ tokens. Request completion/tool handoff offers a final
snapshot, including short rounds. Actual use, value changes, capacity
updates, failure and transfer completion also trigger Algorithm~\ref{alg:placement}.
Candidates cover prefix extension, promotion/demotion, valid final-copy
removal and Peer hot-set replacement. A budgeted refresh revisits stale
keys. If an extension cannot commit, $M$ stays unchanged; copy lag may
make $N>Z$.

\par\medskip\noindent\begin{minipage}{\columnwidth}
\raggedright
 \refstepcounter{designalgorithm}\label{alg:placement}
 \textbf{Algorithm \thedesignalgorithm: Online three-tier placement}
 \par\smallskip
 \begin{tabularx}{\columnwidth}{@{}r@{\quad}X@{}}
 \toprule
 1 & \textsc{Reconcile}$(e)$; \textsc{SettleLowerResidents}$(e)$ \\
 2 & \textbf{if} \textsc{NewUse}$(e)$ \textbf{then} \textsc{RenewAgent}$(e)$ \\
 3 & $S\gets\textsc{RevalueAndSnapshot}(e)$ \\
 4 & $h\gets\textsc{ShrinkRequired}(S)$; $(P^*,d^*)\gets(\varnothing,\infty)$ \\
 5 & \textbf{for} $c\in\textsc{Candidates}(S,h)$ \textbf{while} \textsc{WorkLeft}() \\
 6 & \quad \textbf{for} $P\in\textsc{LegalPlans}(S,c,h)$ \\
 7 & \qquad \textbf{if} $\neg h\land\neg\textsc{LowerRentOK}(P)$ \textbf{then continue} \\
 8 & \qquad $d\gets\Delta J(P;S)$ \\
 9 & \qquad \textbf{if} $d<d^*$ \textbf{then} $(P^*,d^*)\gets(P,d)$ \\
 10 & \textbf{if} $P^*=\varnothing$ \textbf{then} \textsc{Defer}$(e)$; \textbf{return} \\
 11 & \textbf{if} $\neg h\land d^*\ge-\epsilon$ \textbf{then return} \\
 12 & \textbf{if not} \textsc{ReserveIfCurrent}$(S,P^*)$ \textbf{then} \\
 13 & \quad \textsc{Requeue}$(e)$; \textbf{return} \\
 14 & \textsc{Submit}$(P^*)$ \quad\textit{completion/failure emits a new event} \\
 \bottomrule
 \end{tabularx}
\end{minipage}\par\medskip

Connector copies each object, confirms the destination, redirects its
location, then releases the source after readers drain. Both physical
copies consume capacity during this handoff. Versions, leases and pool
clocks are rechecked before reservation; rejected previews change neither
rent nor ownership. Trial rent becomes committed at the corresponding
capacity handoff, and admitted objects keep their residual trial credit.
An aborted copy before handoff preserves the old location; later failure
cannot undo completed moves or rent. Only complete dependency coverage advances the
checkpoint. This bounded search offers no optimality or inherited
competitive guarantee.

\subsection{Adapter: Guarded AIMD for Peer Capacity}
\label{sec:adapter}

Master selects content within a quota; Adapter adjusts that quota as
foreground inference demand changes. It uses pressure-driven multiplicative
decrease and stable-window additive increase~\cite{chiu1989aimd}.
The total Peer authorization has a physical ceiling
\begin{equation}
 C=\max\{0,\min(B_{\mathrm{cap}},
 B_{\mathrm{pool}}-B_{\mathrm{serve}}-B_{\mathrm{reserve}})\}.
 \label{eq:peer-ceiling}
\end{equation}
The pool excludes weights/fixed allocations; serving high watermark
excludes Peer, and reserve covers growth/workspace. Existing Peer bytes
are not subtracted twice. $B_{\rm cap}$ can be a configured pool fraction.

The current quota $B$ and stable-window count $z$ start at zero.
Pressure (allocation stalls, reserve violations or
memory-related admission waits) multiplies $B$ by $0<\gamma<1$. After $m$
stable windows below a lower wait threshold, deferred positive-value Peer
demand, assessed before quota filtering, permits an additive byte step
$\Delta$. Intermediate windows reset
$z$; missing telemetry disables growth. The ceiling clips every update;
emergencies may decrease within a window, but growth requires completed
stable windows. Parameters and window length are configurable.

\par\medskip\noindent\begin{minipage}{\columnwidth}
\raggedright
 \refstepcounter{designalgorithm}\label{alg:adapter}
 \textbf{Algorithm \thedesignalgorithm: Windowed Peer quota}
 \par\smallskip
 \begin{tabularx}{\columnwidth}{@{}r@{\quad}X@{}}
 \toprule
 1 & $C\gets\textsc{SafeCeiling}();\quad B'\gets B$ \\
 2 & \textbf{if} \textsc{Pressure}() \textbf{then} $(B',z)\gets(\gamma B,0)$ \\
 3 & \textbf{else if} \textsc{Stable}() \textbf{then} \\
 4 & \quad $z\gets z+1$ \\
 5 & \quad \textbf{if} $z\ge m$ \textbf{then} \\
 6 & \qquad \textbf{if} \textsc{PeerDemand}() \textbf{then} $B'\gets B+\Delta$ \\
 7 & \qquad $z\gets0$ \\
 8 & \textbf{else} $z\gets0$ \\
 9 & $B\gets\textsc{RoundDownPage}(\min(C,B'))$ \\
 10 & \textsc{PublishBudget}$(B)$; \textsc{EmitCapacityEvent}() \\
 \bottomrule
 \end{tabularx}
\end{minipage}\par\medskip

Adapter sets the total budget; Master chooses $L/K$ at fixed $Z$ and
other placement events. If occupancy plus reservations exceeds a reduced
quota, Master pauses Peer admission and selects objects to demote. Their
DRAM destinations may first require storage demotion. Serving can reuse
the memory only after these transfers and ongoing readers finish; the
new quota is an authorization target, not an instantaneous free-space claim.


\section{Evaluation}
\label{sec:evaluation}

Our evaluation asks five questions: does workflow-aware placement improve
completion after one failure; which recovery and dependency mechanisms
produce that improvement; does it persist under recorded outage timing;
which online mechanisms justify their complexity; and what healthy-service,
resource, and correctness costs accompany it? We evaluate the complete
claim chain, from retained bytes to ready KV to branch completion to
workflow completion, rather than using cache hit rate as its endpoint.

\paragraph{Evidence status.}
This expanded research draft contains reproducible analytical results,
descriptions of public traces, trace-driven model results, and a detailed
physical experiment specification. It does not report an implemented
end-to-end failure-recovery campaign. Each figure identifies its evidence
level. Measured Motivation microbenchmarks retain their own scope;
their timings are not silently substituted for the generic rates below.
The existing DRAM/storage placement results use an earlier whole-checkpoint,
Peer-disabled policy; they do not evaluate the exclusive three-tier
controller. The isolated layer model studies partial DRAM payload
residency and I/O overlap, not the integrated policy. The Design pipeline
reuses existing MiniCPM4-8B calibration with declared transport assumptions;
its Peer case moves selected objects out of DRAM at equal unique coverage.
The static value, isolated rental and layer models test distinct mechanisms
and are not a full-controller run.

\FloatBarrier
\subsection{Experimental Contract and Baselines}
\label{sec:eval-method}

\paragraph{Deployment.}
The primary physical experiment uses failure-separated serving replicas
and independently reachable backing storage. We specify DP=2/4/8 and TP=1/2
as separate cells; a TP2 host-failure test requires at least two compatible
two-GPU replicas on separate hosts. GPU-process, whole-replica, host, and
storage-host failures are distinct experiments. Report surviving capacity
and execution/staging HBM, rather than assuming that a failover target has
unlimited memory. Local NVMe, remote NVMe, and shared filesystem backends
are separate configurations. Same-host SSD does not protect against host
loss. The physical campaign must pin model weights, attention kernels,
engine and controller commits, TP layout, block size, drivers, network
topology, filesystem configuration, and fault-injector version.

\paragraph{Workloads and load.}
Use executable coding-agent and FanOutQA workflows, together with controlled
sequential, balanced fork/join, and heterogeneous fork/join DAGs. Retain
the original graph once, with identical workflow IDs and arrival order;
do not manufacture normal reuse by repeating the DAG. Start with 30
distinct workflows per family and publish their IDs and input hashes.
Replay frozen tool responses for the primary comparison; live tool calls
form a separate validity check. Use supported dense and native-sparse
model paths, including Qwen3-8B and Qwen2.5-7B-Instruct-1M, and contexts
32K/128K/512K/1.01M where executable. Unsupported cells remain missing
with reasons. Arrival traces use 30/60/85\% of one fixed baseline's
sustainable healthy throughput, with identical absolute arrivals across
policies. Report the resulting load after losing a replica.

\paragraph{Equal-resource baselines.}
Recompute keeps durable accepted input/output logs but no additional
fault-safe KV history. Storage-only keeps complete checkpoints in the
capacity tier. Hierarchical LRU and reuse-only share AegisKV's transport,
manifest, safety, and checkpoint-frontier implementation. LRU is the main
same-substrate baseline; reuse-only isolates the benefit of dependency and
failure information. All policies pay for committed, reserved, retiring,
metadata, and staging bytes, and for checkpoint and migration traffic.
The matched online baselines use the same exclusive ownership rule.
Shared objects consume capacity once in steady state; both physical
copies are charged during migration. All-DRAM relaxes
the capacity constraint and is a reference, not an equal-cost competitor.
An exhaustive feasible oracle is restricted to small candidate sets and
the same objective; it is not an online-system oracle.

External comparisons should use pinned D\'ej\`aVu, GhostServe, and LUMEN
artifacts when compatible, and a tuned Mooncake/SGLang or Strata-style
hierarchy~\cite{strati2024dejavu,jayakody2026ghostserve,cao2026lumen,qin2025mooncake,xie2026strata}.
Count replication/parity capacity, coding work, write traffic, and achieved
frontier lag. A reproduced mechanism is labeled as a reimplementation;
none of the present curves is a fabricated result for these named systems.
TOPAS/KVFlow comparisons require the same failure-safe backing substrate
before recovery differences can be attributed to workflow policy.

\paragraph{Outcomes and cohorts.}
For policy $p$ and workflow $w$, the paired failure penalty is
\begin{equation}
 D_{w,p}=\mathrm{WCT}^{\mathrm{fault}}_{w,p}
          -\mathrm{WCT}^{\mathrm{healthy}}_{w,p}.
 \label{eq:eval-penalty}
\end{equation}
Every policy has its own healthy mate. Positive
$D_{w,\mathrm{LRU}}-D_{w,p}$ favors $p$, but we also report absolute healthy
and faulty WCT, completion, and SLO goodput. Separate directly interrupted
workflows, other workflows sharing recovery resources, and an isolated
control cohort. The second cohort establishes interference only when
the actual shared link, queue, or GPU is identified. Report p50/p95/p99
only with adequate completed samples, along with timeouts and the common
observation horizon. Hit rate, restore latency, and workflow latency have
different denominators and are never interchangeable.

\paragraph{Declared analytical instance.}
Four branches A0--A3 have 4/8/12/16\,GiB state, remaining durations
12/10/5/1\,s, and next-use weights 0.15/0.35/0.60/0.85. Rates are
20\,GiB/s DRAM, 4\,GiB/s storage, 12\,GiB/s shared network, and
20\,GiB/s H2D per target, with two surviving targets. The common recovery
floor is 0.25\,s, replay is 0.4\,s, and reconstruction costs 8\,s/GiB.
DRAM is limited to 16\,GiB and storage to 160\,GiB. The event calendar
schedules 0.25\,GiB chunks in round-robin request order through shared
readers, a shared network, and target-local H2D/compute. Tail replay starts
after the required state arrives. Transport overlaps, but this particular
model has no layer-wise replay overlap or online revaluation controller.
Static admission greedily maximizes marginal value per byte with
failure-episode probability 0.05 and next-use probability 0.02 over a
60\,s horizon. These are sensitivity inputs, not observed production rates.

\FloatBarrier
\subsection{One Failure: End-to-End Benefit}
\label{sec:eval-workflow}

\paragraph{Controlled intervention.}
The primary campaign kills the same source replica when a specified
workflow, Agent, and request reaches 80\% of its frozen logical token progress.
The cohort includes every interrupted request on that source, not only
the request that triggers injection. Baseline and candidate runs preserve
the same input, token history, arrival trace, target resources, and failure
identity. Separate the detection floor, routing/admission, checkpoint
lookup, restore, replay, and resumed execution in raw logs. A common
wall-clock kill would confound progress with policy speed.
Mechanism-attribution runs freeze the pre-fault source assignment and
committed coverage. Full-system runs allow policies to change routing
and achieved frontiers, and use the fixed original arrival cohort for
intent-to-treat comparisons; they also report changes in which workflows
were directly interrupted. These two experiments answer different questions.

The analytical grid spans state-size factors 0.5/1/2, slack factors
0/0.3/1/3, and storage rates 1/4/12\,GiB/s, giving 36 configurations.
DRAM scales with state size in this grid. Figure~\ref{fig:eval-workflow-outcomes} shows the distributions. We also compute matched differences for every configuration. Workflow value improves 17 configurations and ties 19;
the absence of a loss in this grid is not a worst-case guarantee.
LRU and reuse-only happen to choose the same order in this declared
instance, so their agreement is not independent evidence against all
reuse policies. Physical comparisons must include other recency/reuse
orders and each workload family separately.

\begin{figure}[!htbp]
 \centering
 \includegraphics[width=\columnwidth]{figures/eval-workflow-outcomes.pdf}
 \caption{Single-failure penalties over 36 declared configurations of state size, residual-branch slack, and storage bandwidth. The same configuration is used for every policy. The CDF weights configurations equally; it is not a distribution of measured workflows. The allocation oracle searches the same whole-group candidate family.}
 \label{fig:eval-workflow-outcomes}
 \Description{Single-failure penalties over 36 declared configurations of state size, residual-branch slack, and storage bandwidth. The same configuration is used for every policy. The CDF weights configurations equally; it is not a distribution of measured workflows. The allocation oracle searches the same whole-group candidate family.}
\end{figure}

In the base case, storage-only adds 5.5583\,s to workflow completion,
LRU adds 3.6208\,s, and workflow value adds 1.6125\,s. The reduction
relative to LRU is 2.0083\,s, or 55.5\% of the \emph{failure penalty};
it is not a 55.5\% end-to-end WCT reduction. The healthy residual join
is 12\,s, giving faulty residual completion times 15.6208 and 13.6125\,s.
Reconstruction's 190.65\,s penalty illustrates the value of reliable
coverage, while the LRU comparison isolates fast-tier placement beyond
that coverage. Measure full workflow lifetimes in the physical experiment;
do not replace them with this residual-window example.

\FloatBarrier
\subsection{One Failure: Explaining the Gain}

\paragraph{Dependency exposure.}
For a single fork/join instance with residual durations $d_i$ and restored
ready times $R_i$, the exact local join penalty is
\begin{equation}
 D=\max_i(d_i+R_i)-\max_i d_i.
\end{equation}
The base join slack is 0/2/7/11\,s. LRU retains the 16\,GiB A3 history;
workflow value retains A0 and A1 (12\,GiB total). Figure~\ref{fig:eval-branch-exposure}
shows why these choices differ. A3's long restoration can remain inside
its available margin, whereas A0 cannot hide any delay. Whole-group
granularity leaves 4\,GiB unused; it is part of the candidate restriction,
not free capacity secretly assigned to another policy.

\begin{figure}[!htbp]
 \centering
 \includegraphics[width=\columnwidth]{figures/eval-branch-exposure.pdf}
 \caption{Where the base-case benefit originates. A0--A3 have join slack 0, 2, 7, and 11\,s. Bars compare that slack with actual modeled restoration completion after the common fault. Speeding up the largest history A3 has little workflow value here; delaying A0 directly extends the join.}
 \label{fig:eval-branch-exposure}
 \Description{Where the base-case benefit originates. A0--A3 have join slack 0, 2, 7, and 11\,s. Bars compare that slack with actual modeled restoration completion after the common fault. Speeding up the largest history A3 has little workflow value here; delaying A0 directly extends the join.}
\end{figure}

\paragraph{Critical-path accounting.}
An ordered counterfactual attributes the final join penalty using controlled interventions. Placement is frozen. We start with negligible
read/network/H2D service time and zero floor/replay, then restore detection-floor, read, network, H2D, and replay costs in that order. Consecutive join differences sum to the final
penalty. This decomposition can assign interactions to different stages
if the intervention order changes. It is therefore an explanatory model
decomposition, not unique causal attribution. In hardware, additionally
record a dependency-linked critical path from event timestamps: sum only
non-overlapping intervals on that path, and report total GPU/I/O work
separately. Neither method may add all SSD, H2D, and GPU durations and
call that wall-clock recovery time.

\paragraph{Cost of prioritization.}
Resource accounting shows the cost of prioritization.
Workflow value reads 28\,GiB from storage, versus LRU's 24\,GiB, and
occupies the complete recovery pool for 7.6833 rather than 6.6833\,s.
It finishes this workflow earlier while consuming the resource longer.
Collect useful bytes, unused prefetched bytes, duplicate reads, checkpoint
writes, and reconstruction work in addition to the critical path. A claim
of lower storage traffic, lower GPU cost, or lower recovery makespan
would be false for this example.

\FloatBarrier
\subsection{One Failure: Collateral Delay and Recovery Backlog}

\paragraph{Shared-resource experiment.}
Cross 1/2/4/8 interrupted workflows with 30/60/85\% pre-fault load and
one/two/four compatible recovery targets. Keep the faulted logical state
and healthy arrival stream fixed. Report direct-workflow delay, collateral
workflow delay, retry-queue occupancy, storage/NIC backlog, and time to
return to the predeclared latency band. Separately remove the common
reader/link by giving targets independent channels. If the collateral
effect disappears, that supports the resource-sharing explanation.
If there is no physical shared bottleneck, do not claim a cross-target
cascade from a model that imposes one.

\begin{figure}[!htbp]
 \centering
 \includegraphics[width=\columnwidth]{figures/eval-burst-interference.pdf}
 \caption{FIFO recovery-pool stress model: a simultaneous burst of identical interrupted workflows shares one two-target pool. The pool is reserved until each workflow's last restore completes. Queue delay is added to the join penalty. This deliberately restrictive policy exposes externalities that an isolated-fault comparison misses.}
 \label{fig:eval-burst-interference}
 \Description{FIFO recovery-pool stress model: a simultaneous burst of identical interrupted workflows shares one two-target pool. The pool is reserved until each workflow's last restore completes. Queue delay is added to the join penalty. This deliberately restrictive policy exposes externalities that an isolated-fault comparison misses.}
\end{figure}

The FIFO stress model in Figure~\ref{fig:eval-burst-interference} holds
one recovery pool until all four branch restores finish. For $n$
simultaneous identical workflows, mean penalty is
$D+(n-1)S/2$, where $S=\max_i R_i$ is pool occupancy. Workflow value
loses its isolated advantage once the additional pool waiting dominates:
at $n=8$, mean penalties are 28.5042\,s for workflow value and
27.0125\,s for LRU. This is a falsification of extrapolating the
isolated placement result to arbitrary bursts, not a measurement of the
proposed load-aware router. With normal I/O using 70\% of the 4\,GiB/s reader, a fluid backlog model exposes a longer drain interval for the larger recovery byte volume. This calculation concerns I/O backlog, not measured TTFT or collateral WCT. The physical test must include concurrent scheduling and
the joint value of affected and collateral workflows.

\FloatBarrier
\subsection{Failure Position and Checkpoint Freshness}

Repeat the same-ID experiment at 20/50/80/95\% logical progress, during
prefill, during decode, during a tool wait, and immediately around a
demotion commit. State precisely whether progress counts accepted total
tokens or generated tokens; never mix denominators. An analytical sweep scales checkpoint bytes with logical progress while holding remaining work fixed.
Real phase sweeps also change the remaining graph, native demand, and
checkpoint frontier, which must be taken from the actual run.

For $Z=64/256/1024$, enumerate every ideal checkpoint phase when presenting
an analytical average, pairing policies at the same $Z$ and phase.
The physical experiment records trigger, copy start, committed frontier,
acknowledged frontier, and retry start. Report lost/uncommitted tail
length, replay time, failed commits, and write backlog. Actual tail length
can exceed $Z$ under backpressure. A short trigger interval can worsen
recovery if it congests the write path. Sweep read/write interference,
write rate, and queue capacity while holding total retention budgets fixed.

\FloatBarrier
\subsection{Real Traces: Provenance and Event Semantics}

\paragraph{Historical node outages.}
LANL documents node interruptions with start/end times, node/system IDs,
and cause categories~\cite{lanlFailureData,schroeder2006failures}. We use
a revision-pinned public redistribution because the original FTP archive
was unavailable during preparation. Its CSV has 23,739 records, whereas
the official page describes 23,741; byte identity and the two-record
difference remain unresolved. This limitation is recorded in the artifact
and must be resolved before presenting the copy as an authenticated
original release. The present result is explicitly for the supplied
redistributed CSV, not for an asserted identical original archive.

We retain rows with purpose exactly \texttt{compute}, exclude 1,648 other
rows and 122 zero-duration rows, and merge overlapping or touching
intervals per node, removing 231 additional interval boundaries. No
negative-duration or unparseable compute intervals occur. The resulting
21,738 intervals cover 3,542 observed nodes across 22 systems. We retain
all recorded cause categories, including interruptions that may involve
maintenance; this is not a hardware-only crash sample. Original row IDs
and every exclusion/merge are recoverable from the cleaning output.

Recorded starts have minute resolution. Equal timestamps identify
coincident \emph{recorded} starts, not proof of a common physical cause.
There are 19,561 within-system start groups, 568 with multiple nodes,
and a maximum of 45 nodes. Systems have separate observation calendars
and are never superimposed as one cluster. These historical CPU/HPC
outages test clustered timing; they do not estimate a 2027 GPU fleet's
failure rate, repair latency, or workload mix. Missing exposure information
also prevents using the observed failing-node count as the denominator
of a fleet-wide availability estimate.

\paragraph{Modern GPU alerts as a separate check.}
The Reflex release contains 98 LLM-training jobs, 344,175 alerts, and
144 faulty node/job labels~\cite{ding2026reflex,reflexDataset}.
We audit these labels separately from outage calendars.
Of the faulty labels, 108 have no alert on the labeled node. Moreover,
timestamps are shifted independently per job and repair times are absent.
We therefore neither stitch jobs into a calendar nor label the first
alert as the fault time. This dataset can support within-job alert-pattern
or diagnosis studies, but it cannot supply an unqualified fail-stop/repair
schedule for this paper. A modern serving outage trace with explicit
event semantics remains an important external-validity target.

\FloatBarrier
\subsection{Trace-Driven Recovery: Benefits and Tail Regressions}

\paragraph{Reproducible conditional replay.}
Each cleaned LANL node-outage start releases one assumed four-branch
workflow requiring recovery. Each original system has an independent
FIFO recovery pool with two surviving targets. Its service occupation
$S_p$ and intrinsic join penalty $D_p$ are the base-case values already
derived from bytes and rates. For consecutive arrivals at $a_i$, pool
waiting is $q_i=\max(0,q_{i-1}+S_p-(a_i-a_{i-1}))$, and workflow
penalty is $q_i+D_p$. Ties use node-ID order. Starts retain their original
spacing. There is no service-time fitting to the trace and no tuned
policy-specific arrival schedule.

This replay is conditional on available independent recovery capacity
and complete committed histories. It does not replay healthy serving,
provision failed nodes, model finite spare exhaustion, or infer which
historical nodes actually hosted active agents. End times are used for
outage characterization, not as measured GPU repair times. Accordingly,
the result is a recovery-load experiment, not an availability or SLO-goodput
measurement of an operating cluster.

\begin{figure}[!htbp]
 \centering
 \includegraphics[width=\columnwidth]{figures/eval-trace-penalty.pdf}
 \caption{LANL-start-time-driven FIFO recovery model, one assumed four-branch workflow per cleaned node outage, preserving per-system start times at their original rate. Curves show 21,738 conditional recovery episodes. The observed input is the outage calendar; recovery durations and workflow penalties are modeled.}
 \label{fig:eval-trace-penalty}
 \Description{LANL-start-time-driven FIFO recovery model, one assumed four-branch workflow per cleaned node outage, preserving per-system start times at their original rate. Curves show 21,738 conditional recovery episodes. The observed input is the outage calendar; recovery durations and workflow penalties are modeled.}
\end{figure}

Figure~\ref{fig:eval-trace-penalty} reports the conditional recovery distribution.
At original spacing, mean modeled penalty is 10.4533\,s for LRU and
9.4673\,s for workflow value, a 9.4\% reduction. Mean queue wait rises
from 6.8324 to 7.8548\,s. The p95 penalty worsens from 37.0375 to
40.0292\,s, and p99 from 204.1208 to 232.1125\,s. Thus this input
supports a conditional average benefit and simultaneously refutes a
blanket tail-latency claim. Per-system aggregation also retains losses, including the negative mean benefit for System 2. Neither thousands of events nor small
arithmetic error bars turn one historical trace into independent serving
experiments.

\paragraph{Controls and full-system extension.}
The evenly spaced control preserves each system's event count and
calendar span while removing original clustering. Timestamp compression
and within-minute spreading distinguish load and timestamp-resolution
sensitivity. Spreading equal-minute starts over 59\,s gives LRU/workflow-value
means 8.0997/7.0224\,s and p99 values 163.8936/191.8852\,s; the adverse
tail remains. These constructed timestamps are not inferred ground truth.
Compression
factors 10/100/1000 increase load explicitly; factor 1 alone retains the
original rate. For a full serving replay, freeze a workload-to-failure-domain
mapping before comparing policies, publish its seed, and use original
start/end intervals to remove and return capacity. Replay arrivals even
during outages, allow failures during unfinished recovery, and retain
incomplete workflows at a common horizon. Repeat mappings at fixed
1/2/4 active workflows per source, different spare fractions, and both
independent and shared backing domains. Preserve correlated event groups
when mapping to fewer machines; never silently serialize a 45-node event
onto one injector. Over-capacity groups must be reported as unsupported
or simulated at the stated scale. Trace amplification and synthetic
repair-time assumptions belong in sensitivity panels, not the primary
real-time trace result.

\FloatBarrier
\subsection{Complete Value and Dependency Ablations}

Every ablation changes one specified mechanism while retaining the same
fault-safe backing, candidate budget, workloads, and physical resources.
The physical matrix includes: no DAG information; independent-request
penalties instead of joint source-event propagation; no normal-reuse
term; no failure term; fixed placement; direct revaluation without rent;
no use-triggered renewal; no current-value refresh; no migration hysteresis;
recency instead of causal
sparse history; serial instead of bounded layer overlap; and locality-only
instead of load-aware recovery routing. Removing a value term never
disables the underlying safety protocol or ordinary cache eligibility.
Correctness mechanisms are tested by injected interleavings, not by
benchmarking an unsafe variant as a valid baseline.

\begin{figure}[!htbp]
 \centering
 \includegraphics[width=\columnwidth]{figures/eval-value-ablation.pdf}
 \caption{Static value ablations over the same 12 size/slack scenarios at 4\,GiB/s storage bandwidth. Each policy has a failure-penalty bar and a normal next-use-cost bar; these are distinct costs, not additive phases. The six policies use identical weights, capacity, and candidate groups. Keeping both outcomes visible exposes the normal-reuse cost of recovery prioritization. This is a static placement model, not the full online system.}
 \label{fig:eval-value-ablation}
 \Description{Static value ablations over the same 12 size/slack scenarios at 4\,GiB/s storage bandwidth. Each policy has a failure-penalty bar and a normal next-use-cost bar; these are distinct costs, not additive phases. The six policies use identical weights, capacity, and candidate groups. Keeping both outcomes visible exposes the normal-reuse cost of recovery prioritization. This is a static placement model, not the full online system.}
\end{figure}

Figure~\ref{fig:eval-value-ablation} reports failure and normal next-use costs together for the same six policies
on the 12-case size/slack subset at 4\,GiB/s storage bandwidth; their
means are not the 36-case bandwidth-mixed means in the outcome CDF.
For each real campaign, additionally log changes in resident objects,
predicted plan value, admitted/rejected migrations, actual critical-path
branches, bytes moved, and controller time. This distinguishes a useful
signal from a coincidentally equivalent placement. The common-failure
example shows why summing isolated losses is insufficient: complementary
required branches can make the marginal value of a bundle different
from the sum of its singleton values. Compare singleton-only candidates
with bounded pairs, then enumerate the oracle on small cases to quantify
candidate-family and search regret separately.

Use factorial experiments for interacting mechanisms: DAG information
$\times$ joint-fault modeling; sparse ranking $\times$ layer overlap;
rental screening $\times$ hysteresis; and load-aware routing $\times$ shared-link
contention. Keep all four cells of each $2\!\times\!2$ design. The
interaction is the paired difference of differences, not the sum of
individual speedups. A tied or negative ablation is informative and
cannot be removed merely because it weakens the proposed narrative.

\FloatBarrier
\subsection{Sparse Demand and Layer-Overlap Ablations}

\paragraph{Controlled mechanism model.}
An eight-layer example uses an SSD payload cost of 4\,ms, H2D cost of
1\,ms, and required compute of 3\,ms/layer at full selected demand.
DRAM retains two payloads, either early or late; transient workspace
holds only the current and next layer. Serial and pipeline modes use
the same demand and resources. Payload density scales I/O for all
policies equally and is not a change to the native attention algorithm.
Figure~\ref{fig:eval-layer-factorial} isolates the placement/scheduling interaction. A separate prefetch-correctness sweep holds required demand fixed to isolate prediction errors.

\begin{figure}[!htbp]
 \centering
 \includegraphics[width=\columnwidth]{figures/eval-layer-factorial.pdf}
 \caption{Equal-budget $2\!\times\!2$ layer-placement/scheduling experiment on eight synthetic layer payloads. Early/late policies retain two payloads; serial/pipeline schedules use the same selected demand, rates, and two-layer transient workspace. The x-axis changes selected payload density for all policies equally, not attention semantics.}
 \label{fig:eval-layer-factorial}
 \Description{Equal-budget $2\!\times\!2$ layer-placement/scheduling experiment on eight synthetic layer payloads. Early/late policies retain two payloads; serial/pipeline schedules use the same selected demand, rates, and two-layer transient workspace. The x-axis changes selected payload density for all policies equally, not attention semantics.}
\end{figure}

At full demand, both serial placements take 56\,ms; pipelining gives
34\,ms with late placement and 32\,ms with early placement. Required
compute is always 24\,ms. The remaining 10 versus 8\,ms is exposed
waiting, not avoided historical prefill. Prediction errors consume
extra reads and force demand-reveal stalls; perfect prediction is only
one sensitivity point.

\paragraph{Matched three-tier extension.}
Hold the same cross-tier checkpoint coverage, fault IDs, native demand
and recovery scheduler fixed; compare Peer disabled, equal-byte
recency/frequency Peer, fixed-depth Peer and per-request L/K revaluation.
Run Peer position $\times$ overlap as a four-cell comparison. Charge
Peer against the same physical HBM pool, including displaced healthy
execution, cross-DP traffic, actual Peer-to-DRAM demotion and temporary
copies. Compare exclusive ownership with duplicated backing separately
at equal physical budgets. Sweep history $M$, actual tail $N$, configured
$Z$, hot-set drift and target rebindings. Report WCT gain, moved bytes and
healthy throughput together. The historical HBM-reuse microbenchmark is
a local-resident reference, not a measurement of Peer placement. These
matched cells remain an experiment specification.

\paragraph{Native implementation experiment.}
Capture the actual per-layer demand union over the retry tail, with
selection metadata and reveal timestamps from the pinned native kernel.
Use historical queries only for prediction and a later held-out replay
window for evaluation; split by document/workflow, not individual blocks.
Compare recency, past access frequency, and the proposed causal profile
at equal physical bytes, mandatory sink/window reservations, and full
checkpoint coverage. Cross each with serial and bounded pipeline restore.
Report all-layer coverage, useful/unused prefetch bytes, missed demand,
reveal-to-ready wait, exposed GPU stalls, workspace peak, total replay,
and WCT. A high early-layer hit rate does not establish high total
coverage. Dense demand is a control, not an implementation failure.
Do not reduce native selections, omit auxiliary state, or count an
approximate answer as exact recovery to make an offload policy faster.

\FloatBarrier
\subsection{Rental Diagnosis and Direct-Revaluation Control}
\label{sec:eval-rental}

The proposed controller combines lower-tier renewable credit, direct
Peer revaluation and complete-plan validation. We retain a negative result for
an earlier rental screen as a regression, not as evidence that the new
controller improves WCT. Its isolated additive model has eight slots,
24 equal-byte objects and 512 deterministic candidate opportunities with
smoothly drifting values. All policies receive the same candidate and
next-use streams, movement cost and replacement threshold. The original
screen renews on candidate revisits, not on scored next-use hits, and
checks one minimum-credit victim. A rejected preview leaves rent unchanged.
This model contains neither multi-Agent shares nor prefix dependencies,
three-tier demotions, physical leases or a serving scheduler.

\begin{figure}[!htbp]
 \centering
 \includegraphics[width=\columnwidth]{figures/eval-rental-cost.pdf}
 \caption{Preserved historical rental regression, in abstract delay-cost
 units. The earlier one-victim trial-rent screen migrates slightly less
 but incurs higher miss-plus-movement cost than direct revaluation.
 These curves do not implement the revised three-tier controller.}
 \label{fig:eval-rental-cost}
 \Description{An earlier rental screen has greater total cost than direct
 value, despite fewer migrations. Its negative result is retained.}
\end{figure}

Figure~\ref{fig:eval-rental-cost} reports direct cost 462.96 and original
rental cost 506.39. They make 108 and 106 migrations and each performs
396 scalar replacement checks, not full DAG passes or CPU measurements.
Removing hysteresis increases migrations to 121; removing renewal reduces
them to 77 while raising cost to 526.05. On the original rental trajectory,
104 rejected opportunities have another resident that satisfies the
current-value gain threshold. These are diagnostic observations, not an
additive causal decomposition of the total gap.

Bounded repairs keep that input and its thresholds unchanged. Renewing
on actual scored hits alone increases cost to 512.67. Refreshing current
value without refunding paid rent reduces it to 477.86; also examining
two low-credit victims reduces it to 466.76. The last repair closes
91.26\% of the original gap but remains 0.82\% above direct revaluation.
Across a fixed grid of four phases, three drift periods and three
capacities, it improves on original rent in all 36 cases, but beats direct
value in only nine. These constructed sensitivity cases are neither
independent repetitions nor held-out workload validation. The repair
also uses 747 replacement checks and 4,088 resident-key updates in the
base case; the result does not establish lower controller overhead.

These diagnostics motivate two explicit design choices: preserve paid
rent when value changes and examine alternative legal victim chains.
They do not measure the revised per-Agent, tiered controller. Its matched
comparison must use the same value model, candidate opportunities,
work budget, safety protocol and resources for direct revaluation,
original screening and the revised mechanism. Ablations remove renewal,
value refresh, shared-Agent aggregation or alternative-victim exploration
individually. Replay multi-round tool waits, growing contexts, shared-role
prefixes, phase changes, full lower tiers and Adapter shrink. Log grants,
paid rent, trial/committed clocks, snapshot invalidations, DAG passes,
controller CPU time, movements, achieved frontiers, misses and WCT. The
negative base case remains part of this regression set.

\FloatBarrier
\subsection{Capacity, Storage, and Recovery Coverage}

Sweep DRAM/unique-checkpoint-working-set ratios 0/0.1/0.25/0.5/1 and
storage ratios 0.5/1/2. Count unique physical bytes and record the
time-varying working set. The finite-storage cells must exercise final-copy
deletion and valid fallback frontiers; they cannot silently turn a DRAM
demotion into unlimited backing storage. Cross equal budgets with
context length and source concurrency. Report unsupported memory/kernel
configurations explicitly. Figure~\ref{fig:eval-capacity-frontier} shows capacity sensitivity; the experiment package also retains bandwidth sweeps. The base model assumes sufficient backing and therefore does not validate final-copy eviction performance.

The value objective is not fault-only WCT. At budgets 16 and 28\,GiB,
the feasible oracle's joint score improves from 0.15049 to 0.14383\,s,
while its fault penalty increases from 1.6125 to 1.9458\,s because
normal reuse also benefits from retention and transfers share resources.
The raw capacity table therefore records both objectives and chosen
objects. A nonmonotone fault-only curve is not evidence that extra
capacity worsens the oracle's optimized joint objective.

\begin{figure}[!htbp]
 \centering
 \includegraphics[width=\columnwidth]{figures/eval-capacity-frontier.pdf}
 \caption{Single-failure DRAM-capacity sweep at fixed storage capacity, rates, and logical fault. Value and oracle policies optimize the joint normal-use/failure score; their fault-only penalty need not decrease monotonically. Equal DRAM capacity is required for deployable-policy comparisons.}
 \label{fig:eval-capacity-frontier}
 \Description{Single-failure DRAM-capacity sweep at fixed storage capacity, rates, and logical fault. Value and oracle policies optimize the joint normal-use/failure score; their fault-only penalty need not decrease monotonically. Equal DRAM capacity is required for deployable-policy comparisons.}
\end{figure}

For each real backend, measure read and write bandwidth, request-size
dependence, concurrent reader/writer contention, commit durability, and
end-to-end GPU-ready time. Cold client cache, warm client cache, and
documented server-cache state are separate cells. Direct I/O controls
the client path only; an uncontrollable shared-filesystem server cache
must remain labeled unknown. Read the same payload and use the same
transport schedule across policies. Add a shared-storage slowdown or
loss experiment to distinguish surviving checkpoint coverage from
nominally low-latency placement. A system cannot recover unavailable
state just because its expected value was high.

\FloatBarrier
\subsection{Prediction, Stale Profiles, and Missing Information}
\label{sec:eval-prediction}

Existing diagnostics use 1,443 request-2 agents in 170 CC/Weka workflows,
with whole-agent remaining time as the target. The five workflow-disjoint
folds give MAEs of 438.6, 406.7 and 407.1\,s for a training-fold median,
basic GBDT and GBDT with history. The paired basic-minus-median interval
$[-80.9,4.9]$\,s crosses zero; the six feature-correlation intervals also
cross zero. These results neither establish a reliable mean improvement
nor validate residual-segment Slack. The raw predictions, complete error
tails and correlation plots remain in the experiment artifact.

Train remaining-time prediction on earlier complete workflows, tune on
a disjoint validation split, and evaluate on later workflows. Features
must be visible at the decision landmark. Report raw remaining-time
error, false-positive slack, critical-path ranking, fallback frequency,
and inference/update cost. A correlation or aggregate WAPE does not
determine a per-agent error CDF. Compare the trained model with duration
medians, zero-slack fallback, and an oracle-duration reference.

The present perturbation model alternates the sign across branches;
at $-25\%$, its 4.5583\,s penalty is worse than LRU's 3.6208\,s.
The physical robustness matrix includes independent and correlated
$\pm10/25/50/100\%$ duration error, missing edges, late tool completion,
stale bandwidth profiles, cold starts, model shifts, and failure-probability
miscalibration. Recompute decisions against actual outcomes rather than
adding an error bar to a nominal curve. Include conservative fallback
and compare the cost of lost opportunity with the cost of false slack.
High predictor accuracy alone is insufficient if it does not change
useful placement decisions.

\FloatBarrier
\subsection{Healthy Service and When Protection Pays}

Run the same workflows without faults at each fixed baseline-defined
load. Include first calls, tool waits, and every Agent request in the
denominator. Report healthy WCT, TTFT, inter-token delay, SLO goodput,
GPU-seconds, controller CPU/memory, checkpoint write volume, and per-tier
request/token hits. Distinguish trace-observed prefix sharing from
realized runtime hits. Compare no new checkpoint protection, LRU,
reuse-only, full value, and no-reuse value with identical storage limits.

A separate M/D/1 calculation with constant service time models I/O queueing only; it does not measure TTFT or GPU interference. For the isolated base case, a simple horizon-level
break-even calculation is $p_f\Delta D-C_h$, where $C_h$ includes
healthy-service and migration costs in the same objective units.
Under this model, a fixed cost of 0.05\,s
requires $p_f>0.0249$ under the assumed 2.0083\,s benefit. This threshold
changes with workload, horizon, and queue externalities. Neither the
illustrative 0.05 probability nor the historical outage trace establishes
the production value of $p_f$. The physical study must report conditional
recovery benefit and amortized net benefit separately. Failure rarity
can invalidate a system-wide value claim even when recovery is faster.

\FloatBarrier
\subsection{Correctness, Safety, and Pressure Transitions}

Exercise ordinary Peer admission with running logical pins, request-end
unpinning, self hot-set replacement and Adapter pressure overrides.
Check complete coverage across all three owning tiers during object
migration. A rebound target must retrieve reachable owners or publish a
valid older-checkpoint/replay fallback after sole-owner loss; it cannot
claim old Peer bytes as local. Test Peer-to-DRAM-to-storage cascades, full
destinations, temporary double occupancy and retained paid rent after
tier round trips. A quota reduction must not report uncompleted demotions
as free memory.

Inject failures immediately before/after reservation, partial copy,
checksum completion, storage commit, directory redirection, token
acknowledgement, and final reader drain. Include duplicate completions,
stale epochs, target loss, missing TP shards/selector state, write timeouts,
and full tiers. A published manifest must identify complete compatible
state; PENDING bytes cannot satisfy a restore. Validate durable accepted
tokens, all-layer replay state, next greedy output under the stated
numerical tolerance, and absence of duplicate acknowledged output.
Tool-side effects require application idempotency evidence separately.

The 24-to-16\,GiB pressure example retains physical bytes until 5.5\,s,
despite the earlier target change and copy completion. Reproduce this
trajectory with overlapping leases, reverse pressure, copy failure,
and last-copy eviction. Check physical, reserved, committed, retiring,
and actually free bytes at every transition; a lower target is not free
memory. Exhaustively enumerate the small state machine and run at least
100 randomized interleavings per transition class in the implementation.
Report the number of distinct injected interleavings, failures, and
unresolved cases. Analytical invariants do not replace protocol execution.

\FloatBarrier
\subsection{Scale and Controller Cost}

Sweep DP=2/4/8, DAG sizes 10/100/1,000, live groups
$10^2/10^3/10^4$, and candidate limits 16/64/256. Independently vary
source-domain concurrency and surviving target count. Record p50/p95/p99
decision time, profile calls, graph visits, edge visits, lock duration,
metadata bytes, plan age, and foreground stalls under a fixed CPU budget.
Small-instance oracle comparisons measure solution quality only where
enumeration is feasible; absence of an oracle at scale is not a zero gap.

The analytical operation-cost projection multiplies an assumed node-visit cost by declared passes/episodes. It establishes neither actual control latency nor full
algorithmic complexity once candidate enumeration and shared-object
owners are included. A useful scaling result must show bounded queues,
acceptable stale-decision rates, and maintained workflow outcomes, not
only a fast isolated scorer. Report saturation and failed configurations
rather than omitting them from a log-scale axis.

\FloatBarrier
\subsection{Statistics, Reproducibility, and Execution Order}
\label{sec:eval-plan}

Start physical primary cells with ten independently launched paired
campaigns. Randomize policy order, retain identical workflow/fault IDs,
and log what each seed changes. Within a campaign, compute each policy's
fault-minus-healthy outcome by workflow ID, then the paired difference
against LRU. Aggregate within campaign before bootstrapping campaigns
(10,000 resamples with a fixed analysis seed); resampling individual
tokens or correlated failures does not create independent repetitions.
For repeated trace windows, use non-overlapping calendar blocks and
report both event-weighted and system/window-weighted summaries.
The deterministic model and finite public-trace descriptions shown here
have zero physical repetitions and no fabricated confidence intervals.

Report common-horizon completion and deadline goodput for every input
workflow, including timeouts. Complete-case latency must use the same
eligible set across policies and show exclusions. If using restricted
mean completion time, state the estimator and censoring assumptions;
do not assign a timeout an invented finish time or let success-only
p99 appear to improve as more workflows fail. Release raw manifests,
arrivals, faults, workflow/request outcomes, dependency-linked events,
prediction snapshots, and all table-generation commands with hashes.
Fixed single-column figure canvases and a separate visual specification
allow replacement data to use the same axes, fonts, and layout.

Execute correctness and backend calibration first, then the same-ID
single-failure matrix and its decomposition, then healthy overhead and
the four factorial ablations, followed by continuous real-trace replay
and pressure/scale tests. Resolve the LANL redistribution discrepancy
and obtain a modern outage trace with trustworthy semantics before
making general fleet claims. The present tail and rental regressions
are explicit decision points for refining or simplifying the design.
An empirical claim is supported only after the corresponding matched
campaign exists; completing an experiment specification does not make
its expected advantage a measured result.

\FloatBarrier
\subsection{Reusable Figures and Measurement Replacement}

All evaluation plots use individual single-column canvases with fixed
physical typography, a separate visual specification, and long-form CSV
tables. PDF, SVG, and PNG exports retain stable filenames. The manifest
records source and table hashes, evidence type, units, aggregation command,
and independent-run count. Measured replacement accepts only tables tied
to declared physical raw logs and campaign-level sample accounting;
changing an evidence badge alone is rejected. A partial measured release
replaces only the completed figure cells while other figures retain their
original labels. Data ranges expand to include negative outcomes and
tails. Changing a metric, such as replacing I/O response time with TTFT,
requires its own explicit semantic specification rather than a renamed
axis. The numerical narrative and captions must be revised with new
results; stable plotting configuration does not automatically validate
an earlier analytical interpretation.

\FloatBarrier

\section{Related Work}
\label{sec:related}

\paragraph{Fault-tolerant LLM serving.}
D\'ej\`aVu streams and replicates KV~\cite{strati2024dejavu}; GhostServe
protects it with host-memory erasure-coded parity~\cite{jayakody2026ghostserve}.
RaidServe addresses TP resilience and imbalance through cyclic placement,
hybrid attention and recovery routing~\cite{xu2026raidserve}. LUMEN
coordinates checkpoints, request redistribution and capacity restoration
under reduced load~\cite{cao2026lumen}. Placement comparisons must use the same strong recovery scheduler.

\paragraph{Hierarchical context caching.}
Mooncake pools KV and schedules around locality~\cite{qin2025mooncake};
its SSD offload spills state under memory pressure~\cite{mooncake2026ssdoffload}.
CachedAttention uses hierarchical caching, layer-wise preload and
asynchronous saving~\cite{gao2024cachedattention}. Strata addresses
fragmented I/O, delay hits and cache-aware scheduling~\cite{xie2026strata}.
HCache restores intermediate activations while balancing I/O and
computation~\cite{gao2025hcache}. Our proposed distinction is common-failure workflow value
with safe demotion and explicit final-copy fallback.

\paragraph{Workflow-aware serving and prediction.}
Parrot optimizes application-level execution~\cite{lin2024parrot}; KVFlow
uses future invocation distance for prefix eviction~\cite{pan2025kvflow}.
Pythia and Cortex exploit workflow semantics~\cite{yu2026pythia,pagonas2025cortex}.
TOPAS schedules retention and execution using remaining paths, reuse and
movement cost~\cite{ni2026topas}; adding a critical-path score to caching
is thus insufficient novelty. The intended distinction is recovery after
a common worker failure with stale checkpoints and finite-tier transitions.
Heddle predicts trajectories~\cite{zhang2026heddle}, which does not establish
AegisKV's remaining-wall-time accuracy.

\paragraph{Placement mechanisms.}
Critical-path analysis~\cite{kelley1959cpm} and cost/size-based caching~\cite{young2002landlord}
provide building blocks. AegisKV's values depend on the DAG, shared queues,
and object locations, so Landlord's competitive bound does not transfer. The
small-instance feasible oracle measures approximation loss. AegisKV
uses lower-tier renewable credit for shared prefixes and tier-conditional
value, direct Peer revaluation and a complete-plan gate. Independent ablations against
direct revaluation must establish the benefit of rent, actual-use renewal
and value refresh, alongside sparse ranking and overlap, under matched
resource and control-work budgets.

\section{Conclusion}
\label{sec:conclusion}

AegisKV combines normal reuse and common-failure workflow cost to manage
finite Peer HBM, DRAM and shared storage. Exclusive ownership distributes
complete checkpoints across these tiers; per-request hot sets accelerate
recovery. The online controller combines lower-tier credit with Peer
revaluation, distinguishing migration from final-copy deletion and
preserving committed coverage through each transition. Analytical
scenarios identify potential benefits and limits. Historical outage timing
shows that average recovery improvements can coexist with tail regressions;
independent ablations identify unvalidated online-policy costs. These
prior model results do not validate the exclusive three-tier controller. The detailed
protocol and single-column figures support the physical correctness and performance
experiments still needed to establish system-level gains.

\paragraph{AI assistance.}
Generative AI assisted drafting, experiment planning, and artifact code;
authors retain responsibility for verification.

\bibliographystyle{ACM-Reference-Format}
\bibliography{references}

% BEGIN MOTIVATION MEASUREMENT APPENDIX
\appendix
\section{Motivation Measurement Protocol}
\label{sec:motivation-protocol}

\paragraph{Scope and configuration.}
Figures~\ref{fig:motivation}(\subref{fig:motivation-recompute}) and
\ref{fig:motivation-recovery-exposure} use the same pinned
Qwen2.5-7B-Instruct-1M weights, SGLang 0.5.10rc0, PyTorch 2.9.1/CUDA 12.8,
TP2, BF16 weights/KV, and batch one on two A100 PCIe 40\,GB GPUs.
They use separate allocations of the same GPU type, with one fixed
repeated synthetic text tokenized by the model tokenizer. The traces
supply workflow positions and length proxies, not benchmark prompts.
Our adapter uses the pinned MInference vertical/slash configuration and
DCA position handling, with sparse attention above 32,768 tokens,
including the final one-token EXTEND. The SM80 sparse-kernel build and
model/attention sources are archived with hashes. GQA K/V are shared;
prefix execution groups at most two query heads, and a single query
batches all heads. These are instrumented worker measurements, rather
than throughput measurements of an unmodified serving frontend.

\paragraph{Common endpoint.}
For $C=M+1$, full prefill computes $M$ tokens from empty KV in chunks of
at most 8,192, then executes one token separately. Reuse consumes exactly
the same sparse-generated $M$-token state, including partial pages, and
executes that final token. Synchronized wall time is reduced by maximum
across the TP ranks and ends at model-body completion with all $C$ tokens'
KV resident, before the output vocabulary projection. Model loading,
backup creation, and correctness checks are outside the timer. Detection,
rerouting, serving-queue delay and an injected failure are absent.

\paragraph{Restore paths.}
HBM begins with resident KV. DRAM copies all historical layers from
pinned host buffers before computation. Shared storage opens, reads and
closes per-rank KV files on the compute node's DPC shared filesystem,
then performs the same H2D copy and computation. These stages do not
overlap; there is no compression or selective loading. The 1,010,000-token
case restores 57,917,382,656 logical history bytes across both ranks;
files additionally contain recorded alignment padding. Direct I/O and
page-residency checks exclude Linux file-page hits, but DPC server and
internal caches are uncontrolled. The curve is a shared-storage path
measurement, not cold-device latency or an identified SSD measurement.
KV files are removed after measurement.

\paragraph{Timing samples.}
Figure~\ref{fig:motivation}(\subref{fig:motivation-recompute}) uses one
allocation for 8,192, 65,536, 262,144 and 1,010,000 tokens. HBM/DRAM have
one warmup and three timings per length; full prefill and shared storage
have one timing each. Short forwards and operator checks precede timing;
64K also runs a complete original all-head sparse reference. There is
no extra complete-prefill warmup at 256K/1.01M or storage-read warmup.
Consequently, these full-prefill/storage points describe individual
runs, not stable latency distributions.
Figure~\ref{fig:motivation-recovery-exposure} uses two allocations at
$C=100{,}865$, $105{,}921$, $108{,}865$ and $113{,}665$, with one complete
warmup and three full-prefill timings at each length. The initial allocation completed three lengths; the last length was
resumed after an early termination signal. Its incomplete timings are
retained in the artifact but excluded as a whole. The figure reports
medians and min--max ranges. Repetitions are timing samples of the same input,
not independent workflows or confidence intervals. Original timings and
all samples are retained in the artifact.

\paragraph{Validation and workflow mapping.}
Complete-history hashes, final hidden states, all-layer new KV and greedy
outputs are checked across reconstruction/reuse paths. Repeated full
prefills additionally check output agreement, and CUDA-event timelines
check serialized copies. These checks validate matched sparse-state
execution, not equivalence to dense attention or long-context task quality.
For Figure~\ref{fig:motivation-recovery-exposure}, the four history lengths
are the active requests' trace hash-block counts multiplied by 64.
They are size proxies, not exact original-model token counts or observed
KV frontiers at failure. The example was selected for four parallel
branches and a gap near 100\,s; it is not its workflow's maximum-gap group.
Each counterfactual separately inserts a measured full-prefill duration
as an additional branch pause at 29.5\,s after the fork. The all-branch
join, other execution and queues stay fixed. Thus only reconstruction
cost is measured: the plotted WCT increment is a conditional calculation,
not observed fault recovery in the Claude Code workflow.
% END MOTIVATION MEASUREMENT APPENDIX
\end{document}
